\begin{lemma}
\label{lemma-top-chern-class}
In the situation described just above assume $\dim_\delta(X') = n$,
that $f^*\mathcal{E}$ has constant rank $r$, that
$\dim_\delta(Z(s)) \leq n - r$, and that for every generic point
$\xi \in Z(s)$ with $\delta(\xi) = n - r$ the ideal of $Z(s)$
in $\mathcal{O}_{X', \xi}$ is generated by a regular sequence
of length $r$. Then
$$
c_r(\mathcal{E}) \cap [X']_n = [Z(s)]_{n - r}
$$
in $\CH_*(X')$.
\end{lemma}

\begin{proof}
Since $c_r(\mathcal{E})$ is a bivariant class
(Lemma \ref{lemma-cap-cp-bivariant})
we may assume $X = X'$ and we have to show that
$c_r(\mathcal{E}) \cap [X]_n = [Z(s)]_{n - r}$ in $\CH_{n - r}(X)$.
We will prove the lemma by induction on $r \geq 0$. (The case
$r = 0$ is trivial.) The case $r = 1$
is handled by Lemma \ref{lemma-geometric-cap}. Assume $r > 1$.

\medskip\noindent
Let $\pi : P \to X$ be the projective space bundle associated to
$\mathcal{E}$ and consider the short exact sequence
$$
0 \to \mathcal{E}' \to \pi^*\mathcal{E} \to \mathcal{O}_P(1) \to 0
$$
By the projective space bundle formula
(Lemma \ref{lemma-chow-ring-projective-bundle})
it suffices to prove the equality after pulling back by $\pi$.
Observe that $\pi^{-1}Z(s) = Z(\pi^*s)$ has $\delta$-dimension
$\leq n - 1$ and that the assumption on regular sequences at
generic points of $\delta$-dimension $n - 1$ holds by
flat pullback, see
Algebra, Lemma \ref{algebra-lemma-flat-increases-depth}.
Let $t \in \Gamma(P, \mathcal{O}_P(1))$ be the image of $\pi^*s$.
We claim
$$
[Z(t)]_{n + r - 2} = c_1(\mathcal{O}_P(1)) \cap [P]_{n + r - 1}
$$
Assuming the claim we finish the proof as follows.
The restriction $\pi^*s|_{Z(t)}$ maps to zero in
$\mathcal{O}_P(1)|_{Z(t)}$ hence comes from a unique
element $s' \in \Gamma(Z(t), \mathcal{E}'|_{Z(t)})$.
Note that $Z(s') = Z(\pi^*s)$ as closed subschemes of $P$.
If $\xi \in Z(s')$ is a generic point with $\delta(\xi) = n - 1$,
then the ideal of $Z(s')$ in $\mathcal{O}_{Z(t), \xi}$
can be generated by a regular sequence of length $r - 1$: it is generated by
$r - 1$ elements which are the images of $r - 1$ elements in
$\mathcal{O}_{P, \xi}$ which together with a generator of the
ideal of $Z(t)$ in $\mathcal{O}_{P, \xi}$ form a regular sequence
of length $r$ in $\mathcal{O}_{P, \xi}$. Hence we can apply the
induction hypothesis to $s'$ on $Z(t)$ to get
$c_{r - 1}(\mathcal{E}') \cap [Z(t)]_{n + r - 2} = [Z(s')]_{n - 1}$.
Combining all of the above we obtain
\begin{align*}
c_r(\pi^*\mathcal{E}) \cap [P]_{n + r - 1}
& =
c_{r - 1}(\mathcal{E}') \cap c_1(\mathcal{O}_P(1)) \cap [P]_{n + r - 1} \\
& =
c_{r - 1}(\mathcal{E}') \cap [Z(t)]_{n + r - 2} \\
& =
[Z(s')]_{n - 1} \\
& = [Z(\pi^*s)]_{n - 1}
\end{align*}
which is what we had to show.

\medskip\noindent
Proof of the claim. This will follow from an application of
the already used Lemma \ref{lemma-geometric-cap}.
We have $\pi^{-1}(Z(s)) = Z(\pi^*s) \subset Z(t)$.
On the other hand, for $x \in X$ if $P_x \subset Z(t)$, then
$t|_{P_x} = 0$ which implies that $s$ is zero in the fibre
$\mathcal{E} \otimes \kappa(x)$, which implies $x \in Z(s)$.
It follows that $\dim_\delta(Z(t)) \leq n + (r - 1) - 1$.
Finally, let $\xi \in Z(t)$ be a generic point with
$\delta(\xi) = n + r - 2$. If $\xi$ is not the generic point
of the fibre of $P \to X$ it is immediate that
a local equation of $Z(t)$ is a nonzerodivisor in $\mathcal{O}_{P, \xi}$
(because we can check this on the fibre by
Algebra, Lemma \ref{algebra-lemma-grothendieck}).
If $\xi$ is the generic point of a fibre, then $x = \pi(\xi) \in Z(s)$
and $\delta(x) = n + r - 2 - (r - 1) = n - 1$. This is a contradiction
with $\dim_\delta(Z(s)) \leq n - r$ because $r > 1$
so this case doesn't happen.
\end{proof}